\documentclass[11pt,letterpaper]{article}

\usepackage[margin=0.65in]{geometry}
\usepackage{amsmath,amssymb}
\usepackage{booktabs}
\usepackage{array}
\usepackage{tabularx}
\usepackage{graphicx}
\usepackage{float}
\usepackage{caption}
\usepackage{subcaption}
\usepackage[hidelinks]{hyperref}
\usepackage{parskip}
\usepackage{microtype}
\usepackage[authoryear,round]{natbib}
\usepackage{times}
\usepackage{titlesec}
\usepackage[table]{xcolor}
\usepackage{enumitem}
\usepackage{multirow}

% ── Colors ──────────────────────────────────────────────────────────────────────
\definecolor{navy}{HTML}{1F4E79}
\definecolor{midblue}{HTML}{2E75B6}
\definecolor{rowblue}{HTML}{EBF3FB}
\definecolor{headblue}{HTML}{D5E8F0}
\definecolor{rowgray}{gray}{0.94}

% ── Section formatting ───────────────────────────────────────────────────────────
\titleformat{\section}
  {\large\bfseries\color{navy}}{\thesection.}{0.5em}{}
  [\vspace{-4pt}{\color{midblue}\rule{\linewidth}{1.2pt}}\vspace{2pt}]
\titleformat{\subsection}
  {\normalsize\bfseries\color{navy}}{\thesubsection}{0.5em}{}
\titleformat{\paragraph}[runin]
  {\normalsize\bfseries\color{navy}}{}{0pt}{}[.\enspace]
\titlespacing{\section}{0pt}{12pt}{5pt}
\titlespacing{\subsection}{0pt}{8pt}{3pt}
\titlespacing{\paragraph}{0pt}{6pt}{0pt}

% ── Caption formatting ───────────────────────────────────────────────────────────
\captionsetup{font=small,labelfont={bf,color=navy},textfont=it,
              labelsep=period,skip=4pt}
\captionsetup[table]{position=top,skip=4pt}
\captionsetup[subtable]{font=small,labelfont=bf,textfont=it,skip=3pt}

\setlength{\parskip}{5pt}
\setlength{\parindent}{0pt}
\setlength{\extrarowheight}{2pt}

% ── Column types ─────────────────────────────────────────────────────────────────
\newcolumntype{R}{>{\raggedleft\arraybackslash}X}
\newcolumntype{L}{>{\raggedright\arraybackslash}X}

% ── Title ─────────────────────────────────────────────────────────────────────────
\title{%
  \vspace{-1.8em}%
  {\color{navy}\textbf{\large Multi-Asset Risk Parity with Regime-Conditional Tactical Tilts}}\\[5pt]
  \normalsize\mdseries\color{black}
  Group 13\enspace$|$\enspace
  MGMTMFE 431, Quantitative Asset Management\enspace$|$\enspace
  UCLA Anderson, Spring 2026
}
\author{Mayuri Salunke \and Aarya Parashar \and Yixin Xia \and Siyao Pei}
\date{\normalsize May 3, 2026}

\begin{document}
\maketitle
\vspace{-14pt}
{\color{midblue}\rule{\linewidth}{1.5pt}}
\vspace{4pt}

% ══════════════════════════════════════════════════════════════════════════════
\section{Introduction and Motivation}

Classical mean-variance optimization is compelling in theory but unreliable in
practice. Estimation error in expected returns produces extreme, unstable weights
that systematically underperform naive alternatives \citep{DeMiguel2009,Michaud1989}.
The standard remedy, shrinkage or constraints on $\hat{\mu}$, attenuates the
problem without eliminating it. A cleaner solution is to build the \emph{base}
portfolio on covariance structure alone, then layer a disciplined, forecast-driven
tactical overlay on top. This two-layer design separates robust risk budgeting from
return timing, allowing each layer to do what it does best.

Our strategy has four interconnected motivations:

\begin{itemize}[leftmargin=1.4em,itemsep=2pt,topsep=4pt]
  \item \textbf{Risk parity over mean-variance.} Equal Risk Contribution (ERC)
    requires no return forecasts and allocates capital so each asset contributes
    equally to total portfolio variance \citep{Maillard2010}. It outperforms naive
    $1/N$ on a drawdown-adjusted basis in multi-asset settings where volatility
    differs substantially across asset classes.
  \item \textbf{Genuine diversification.} We construct a seven-asset universe
    spanning equities, government bonds, credit, commodities, and a long/short
    commodity carry strategy. The final asset, the AQR commodity carry L/S return,
    has near-zero correlation with equity ($\rho = 0.035$) and with the long-only
    commodity return ($\rho = -0.04$), making it a genuine source of alpha
    orthogonal to the standard eight-factor risk model.
  \item \textbf{Regime-conditional forecasting.} Return predictability is
    state-dependent: macro predictors are more informative during financial stress
    \citep{Adrian2019,Cenesizoglu2012}. We build a financial conditions index (FCI)
    from credit spreads, stock variance, and the term spread, then run separate
    regressions within each regime to generate combination forecasts.
  \item \textbf{Disciplined overlay.} Raw forecasts are shrunk 70\% toward zero
    \citep{CampbellThompson2008}, passed through an expanding-window confidence
    gate, and scaled by a continuous regime score before tilting. The sign-and-cap
    rule (maximum $\pm 10$ pp per asset).
\end{itemize}

% ══════════════════════════════════════════════════════════════════════════════
\section{Data}

Table~\ref{tab:data} summarizes all data sources. The sample runs from July 1990
through December 2024 (414 months), with the start date set by the availability of
international equity data.

\begin{table}[H]
\small\centering
\caption{Data sources and coverage.}
\label{tab:data}
\setlength{\tabcolsep}{6pt}
\begin{tabularx}{\linewidth}{llLr}
\toprule
\textbf{Series} & \textbf{Source} & \textbf{Description} & \textbf{Coverage} \\
\midrule
\rowcolor{rowblue}
US\_EQ    & Ken French & US equity market excess return (Mkt-RF) & 1990--2024 \\
INTL\_EQ  & Ken French & Developed ex-US equity excess return & 1990--2024 \\
\rowcolor{rowblue}
LONG\_TREAS & Goyal-Welch & Long-term US Treasury excess return (\textit{ltr}) & 1990--2024 \\
CREDIT    & Goyal-Welch & IG corporate bond excess return (\textit{corpr}) & 1990--2024 \\
\rowcolor{rowblue}
COMMOD    & AQR & Equal-weight commodity excess return & 1990--2024 \\
INT\_BOND & FRED GS5   & 5-year Treasury excess return (duration approx.) & 1990--2024 \\
\rowcolor{rowblue}
LS\_COMMOD & AQR & Long/short commodity carry excess return & 1990--2024 \\
\midrule
7 GW predictors & Goyal-Welch & T-bill, term spread, credit spread, div.\ yield, & 1990--2024 \\
 & & inflation, net issuance, stock variance & \\
\rowcolor{rowblue}
ip\_mom & FRED INDPRO & 12-month industrial production growth (2-month lag) & 1990--2024 \\
cpi\_mom & FRED CPI & 12-month CPI inflation rate (2-month lag) & 1990--2024 \\
\rowcolor{rowblue}
slope\_5\_10 & FRED GS10 & 10yr minus 5yr Treasury yield slope & 1990--2024 \\
backwardation & AQR & Aggregate commodity backwardation/contango score & 1990--2024 \\
\bottomrule
\end{tabularx}
\end{table}

\paragraph{The LS\_COMMOD asset}
The long/short commodity carry series from \citet{Levine2018} deserves special
attention. Its standalone Sharpe ratio is 0.60, its correlation with US equity is
0.035, and its correlation with the equal-weight commodity long return is $-0.04$.
This near-orthogonality to all equity and bond factors makes it a genuine diversifier
and the primary source of factor-adjusted alpha in our strategy.

\paragraph{INT\_BOND}
The 5-year Treasury excess return is constructed from FRED's GS5 series using the
standard duration approximation $r_t \approx y_{t-1}/12 - D_{\text{mod}}(y_t - y_{t-1}) - r_f$,
where $D_{\text{mod}} = 4.4$ years. It has Sharpe 0.46, correlation $-0.04$ with
US equity, and correlation 0.56 with the long Treasury, providing meaningful
interest-rate diversification within the fixed-income sleeve.

% ══════════════════════════════════════════════════════════════════════════════
\section{Portfolio Construction}

\subsection{Layer 1: Equal Risk Contribution Base}

Let $\Sigma_t$ denote the covariance matrix estimated at month $t$. The ERC base
weights $\mathbf{w}^{\text{ERC}}_t \in \mathbb{R}^7$ satisfy:
\begin{equation}
  w^{\text{ERC}}_{i,t}\,(\Sigma_t\,\mathbf{w}^{\text{ERC}}_t)_i
  \;=\;
  w^{\text{ERC}}_{j,t}\,(\Sigma_t\,\mathbf{w}^{\text{ERC}}_t)_j
  \quad \forall\, i \ne j,
  \label{eq:erc}
\end{equation}
solved via the convex log-barrier dual of \citet{Maillard2010}: $\min_{\mathbf{y}>0}\,
\frac{1}{2}\mathbf{y}'\Sigma_t\mathbf{y} - \mathbf{1}'\log\mathbf{y}$,
with $\mathbf{w} = \mathbf{y}/\mathbf{1}'\mathbf{y}$.

\paragraph{Covariance estimation}
We use three improvements over a plain sample covariance:
(i)~\textbf{Regime-conditional EWMA}: $\lambda = 0.90$ in turbulent months (more
weight on recent dislocations) and $\lambda = 0.96$ in calm months, interpolated
continuously via the regime score. (ii)~\textbf{Ledoit-Wolf shrinkage} toward a
scaled identity, with the shrinkage coefficient estimated analytically on each
rolling 60-month window \citep{LedoitWolf2004}. (iii)~\textbf{Continuous blending}:
the blend proportion is $\hat{\sigma}_{\text{LW}}$ for the identity component and
$1-\hat{\sigma}_{\text{LW}}$ for the EWMA component.

\subsection{Layer 2: Regime-Conditional Tactical Overlay}

\paragraph{Regime scoring}
We form an ABG-style financial conditions index \citep{Adrian2019}:
\begin{equation}
  \text{FCI}_t \;=\; z(\text{credit spread}_t) + z(\text{svar}_t) - z(\text{term spread}_t),
\end{equation}
where $z(\cdot)$ denotes an expanding $z$-score using only data strictly before $t$.
Rather than a binary calm/turbulent flag, we define a \textbf{continuous regime score}
$s_t = \Pr(\text{FCI}_{<t} < \text{FCI}_t)$, the expanding percentile rank. This
smooths regime transitions, reduces unnecessary turnover, and scales the overlay
continuously: $s_t = 0$ implies no tilt; $s_t = 1$ implies full tilt.

\paragraph{Combination forecasts}
For each asset $i$ and each of 11 predictors $j$ (7 Goyal-Welch, 3 macro, backwardation
score, plus asset-specific TSMOM for 12 signals total per asset) we estimate
\begin{equation}
  r_{i,\tau+1} = \alpha_{i,j,s} + \beta_{i,j,s}\,x_{j,\tau} + \varepsilon_{i,\tau+1}
\end{equation}
separately within the current binary regime (calm/turbulent), using all historical
regime-matched observations. We combine by equal weighting \citep{Rapach2010}:
\begin{equation}
  \hat{r}_{i,t+1}^{\text{combo}} \;=\; \tfrac{1}{12}\textstyle\sum_{j=1}^{12}\hat{r}_{i,j,t+1}.
\end{equation}

\paragraph{Forecast conditioning and overlay rule}
The raw combination forecast is transformed in three steps before entering the tilt:
\begin{enumerate}[leftmargin=1.4em,itemsep=1pt,topsep=3pt]
  \item \textbf{Shrinkage.} $\tilde{\mu}_{i,t} = 0.30 \times \hat{r}_{i,t+1}^{\text{combo}}$.
    Trusting only 30\% of the raw signal robustly reduces estimation noise
    \citep{CampbellThompson2008}.
  \item \textbf{Confidence gate.} Asset $i$ receives a non-zero tilt only if
    $|\tilde{\mu}_{i,t}| > \sigma_{\hat{r}_i}^{\text{expand}}$, where
    $\sigma_{\hat{r}_i}^{\text{expand}}$ is the expanding standard deviation of
    past forecasts (no lookahead). This filters the noisiest signals.
  \item \textbf{Sign-cap tilt, scaled by regime.} Define
    $x_{i,t} = \text{sign}(\tilde{\mu}_{i,t}^{\text{gated}}) \times 0.10$,
    projected to $\mathbf{1}'\mathbf{x}=0$ and clipped to
    $[-0.10,\,0.10]\cap[-w^{\text{ERC}}_{i,t},\,0.10]$.
    The final weight is $w_{i,t} = w^{\text{ERC}}_{i,t} + s_t \cdot x_{i,t}$.
\end{enumerate}

The portfolio is rebalanced monthly. All predictors are lagged one month (two months
for INDPRO and CPI) before use, ensuring strict no-lookahead compliance.

% ══════════════════════════════════════════════════════════════════════════════
\section{Backtesting Results}

\subsection{Setup}

The burn-in period is July 1990 through June 1995. The out-of-sample (OOS)
evaluation period is July 1995 through December 2024 (354 months). Benchmarks: ERC
base (7-asset), 60/40 (60\% US equity + 40\% long Treasury), vol-targeted 60/40
(10\% target, leverage cap $1.5\times$), and equal-weighted $1/N$. Transaction
costs are deducted using period-specific monthly turnover at 10 bps one-way.

\subsection{Headline Performance}

\begin{table}[H]
\small\centering
\caption{Annualized OOS performance metrics, July 1995 -- December 2024 (gross of costs; net figures in Table~\ref{tab:tc}).}
\label{tab:perf}
\setlength{\tabcolsep}{5pt}
\begin{tabularx}{\linewidth}{Lrrrrrrr}
\toprule
\textbf{Portfolio} & \textbf{Mean} & \textbf{Vol} & \textbf{Sharpe} & \textbf{MaxDD} & \textbf{Calmar} & \textbf{Alpha} & \textbf{$t$-stat} \\
 & (\%) & (\%) & & (\%) & & (\%) & \\
\midrule
\rowcolor{rowblue}
ERC + Tactical \textbf{(ours)} & \textbf{4.70} & \textbf{5.32} & \textbf{0.884} & $\mathbf{-12.2}$ & \textbf{0.384} & \textbf{1.344} & \textbf{1.95} \\
ERC base (7 assets)     & 4.54 & 5.23 & 0.867 & $-$12.9 & 0.352 & 1.173 & 1.72 \\
60/40                   & 6.97 & 9.63 & 0.724 & $-$30.7 & 0.227 & 0.000 & --- \\
Vol-Targeted 60/40      & 7.35 & 10.41 & 0.706 & $-$34.6 & 0.212 & $-$0.237 & $-$0.39 \\
Equal-weighted (1/N)    & 5.26 & 6.72 & 0.782 & $-$20.9 & 0.251 & 1.005 & 1.22 \\
\bottomrule
\end{tabularx}
{\footnotesize Alpha is annualized (\%), estimated against the 8-factor model (FF5 + MOM + TERM + DEF) with Newey-West SEs (6 lags). 60/40 alpha is zero by construction (it lies in the factor span of MKT + TERM).}
\end{table}

Our strategy achieves a Sharpe ratio of \textbf{0.884}, the highest of any portfolio
examined, against 0.867 for the ERC base and 0.724 for 60/40. More importantly, it
does so with maximum drawdown of only $-12.2\%$ versus $-30.7\%$ for 60/40 and
$-20.9\%$ for equal-weighted. The Calmar ratio of 0.384 is 69\% above 60/40 (0.227). The stability of our strategy's drawdown profile is best appreciated by noting its peak loss of -12.2\% occurred during the 2008 crisis, a period when the 60/40 benchmark suffered more than double that magnitude.

The 8-factor alpha of \textbf{1.344\% per year} ($t = 1.95$) is economically and
statistically meaningful, and higher than the ERC base (1.173\%, $t = 1.72$).
This improvement directly reflects the addition of LS\_COMMOD, whose returns are
orthogonal to all equity and bond risk factors. The R\textsuperscript{2} of 68.4\%
is deliberately lower than the ERC base (69.8\%) because LS\_COMMOD returns are
not spanned by the factor model, which is exactly the property we want: a genuine
out-of-model return source.

\begin{figure}[H]
\centering
\begin{subfigure}[t]{0.49\linewidth}
  \includegraphics[width=\linewidth]{fig1_cumulative.png}
  \caption{Log cumulative excess returns. Our strategy (blue) compounds steadily with lower volatility than 60/40. Shaded bands: NBER recessions.}
\end{subfigure}
\hfill
\begin{subfigure}[t]{0.49\linewidth}
  \includegraphics[width=\linewidth]{fig2_drawdowns.png}
  \caption{Drawdown series. The strategy (blue, nearly invisible near zero) sustains the shallowest drawdowns across every major stress episode including 2001, 2008, and 2022.}
\end{subfigure}
\caption{Out-of-sample performance, July 1995 -- December 2024.}
\label{fig:perf}
\end{figure}

\subsection{Factor Model Analysis}

\begin{table}[H]
\small\centering
\caption{8-factor alpha regression (FF5 + MOM + TERM + DEF), Newey-West standard errors (6 lags).}
\label{tab:alpha}
\setlength{\tabcolsep}{6pt}
\begin{tabularx}{\linewidth}{Lrrr}
\toprule
\textbf{Portfolio} & \textbf{Alpha (\%)} & \textbf{$t$-stat} & \textbf{$R^2$ (\%)} \\
\midrule
\rowcolor{rowblue}
ERC + Tactical (ours) & 1.344 & 1.95 & 68.4 \\
ERC base (7 assets)   & 1.173 & 1.72 & 69.8 \\
60/40                 & 0.000 & ---  & 100.0 \\
Equal-weighted (1/N)  & 1.005 & 1.22 & 74.2 \\
\bottomrule
\end{tabularx}
\smallskip
\begin{tabularx}{\linewidth}{Lrrrrrrrr}
\toprule
\multicolumn{9}{l}{\textit{Factor loadings, ERC + Tactical:}} \\
\midrule
 & MKT & SMB & HML & RMW & CMA & MOM & TERM & DEF \\
\midrule
\rowcolor{rowblue}
$\hat\beta$ & 0.178 & 0.036 & $-$0.016 & 0.026 & 0.047 & 0.004 & 0.349 & 0.284 \\
\bottomrule
\end{tabularx}
\end{table}

The dominant factor exposures are MKT (0.178) and TERM (0.349), reflecting the
multi-asset nature of the portfolio. DEF loading (0.284) captures the credit and
commodity carry sleeves. Crucially, all equity style factor loadings (SMB, HML,
RMW, CMA, MOM) are near zero, confirming the strategy does not inadvertently
harvest known equity anomalies. The positive and near-significant alpha ($t = 1.95$)
reflects the genuine orthogonality of LS\_COMMOD to this factor space.

The rolling 60-month alpha in Figure~\ref{fig:alpha} confirms the positive alpha was strongest during 2004--2011, when commodity carry premia were elevated, and has compressed since 2018. This likely reflects both increased market efficiency as carry trades became institutionalized and a shift in macro-regime dynamics where the FCI's traditional components (like credit spreads) were dampened by prolonged liquidity injections, reducing the "signal-to-noise" ratio of tactical tilts in the post-QE era.

\begin{figure}[H]
\centering
\begin{subfigure}[t]{0.49\linewidth}
  \includegraphics[width=\linewidth]{fig3_rolling_sharpe.png}
  \caption{Rolling 36-month Sharpe ratio. Strategy (blue) consistently above 60/40 (green) and closely tracks ERC base (orange) but with additional carry return.}
\end{subfigure}
\hfill
\begin{subfigure}[t]{0.49\linewidth}
  \includegraphics[width=\linewidth]{fig5_rolling_alpha.png}
  \caption{Rolling 60-month 8-factor alpha. Positive in 2004--2011 during the commodity supercycle; near zero post-2018 as carry premia compressed.}
\end{subfigure}
\caption{Rolling performance metrics. Shaded bands: NBER recessions.}
\label{fig:alpha}
\end{figure}

\subsection{Forecast Quality}

\begin{table}[H]
\small\centering
\caption{Combination forecast quality: OOS $R^2$ (\%) and directional accuracy (\%) by regime.}
\label{tab:forecast}
\setlength{\tabcolsep}{5pt}
\begin{tabularx}{\linewidth}{Lrrrr}
\toprule
\textbf{Asset} & \textbf{OOS $R^2$ (\%)} & \textbf{DA: All (\%)} & \textbf{DA: Calm (\%)} & \textbf{DA: Turbulent (\%)} \\
\midrule
\rowcolor{rowblue}
US Equity     & $-$1.195 & 63.0 & 67.4 & 60.4 \\
Intl Equity   & $-$2.209 & 54.0 & 54.5 & 53.6 \\
\rowcolor{rowblue}
Long Treasury & $-$0.708 & 50.8 & 55.3 & 48.2 \\
Credit        & $-$0.435 & 56.5 & 62.1 & 53.2 \\
\rowcolor{rowblue}
Commodities   & $-$2.243 & 58.5 & 58.3 & 58.6 \\
INT\_BOND     & $+$0.540 & 53.7 & 49.2 & 55.9 \\
\rowcolor{rowblue}
LS\_COMMOD    & $-$1.577 & 59.3 & 62.1 & 57.7 \\
\bottomrule
\end{tabularx}
{\footnotesize OOS $R^2$: Campbell-Thompson statistic vs.\ expanding historical mean. DA = directional accuracy (fraction of months where sign of forecast matches sign of realized return).}
\end{table}

Negative OOS $R^2$ values are consistent with the equity-predictability literature
\citep{GoyalWelch2008}. However, as \citet{Cenesizoglu2012} show, negative $R^2$ is
compatible with positive portfolio value when directional accuracy exceeds 50\%. While these negative values suggest the historical mean is a tough-to-beat hurdle, the consistency of directional accuracy above 50\% across diverse assets indicates that the tactical overlay captures persistent, albeit small, economic signals rather than mere sampling noise.
Five of seven assets exceed 50\% directional accuracy unconditionally, and US equity
(63.0\%) and credit (56.5\%) show meaningful predictability. INT\_BOND is the only
asset with positive OOS $R^2$ (0.540\%), driven by the predictive power of the
5yr-10yr yield slope for intermediate bond returns.

\subsection{Regime-Conditional Behavior}

\begin{table}[H]
\small\centering
\caption{Annualized mean excess return by regime (\%) and Welch two-sample $t$-test.}
\label{tab:regime}
\setlength{\tabcolsep}{6pt}
\begin{tabularx}{\linewidth}{Lrrrr}
\toprule
\textbf{Portfolio} & \textbf{Calm (\%)} & \textbf{Turbulent (\%)} & \textbf{Diff (\%)} & \textbf{$t$-stat} \\
\midrule
\rowcolor{rowblue}
ERC + Tactical (ours) & 4.83 & 4.62 & 0.21 & 0.11 \\
ERC base (7 assets)   & 4.85 & 4.36 & 0.49 & 0.27 \\
60/40                 & 8.32 & 6.17 & 2.15 & 0.64 \\
Equal-weighted (1/N)  & 5.41 & 3.95 & 1.46 & 0.57 \\
\bottomrule
\end{tabularx}
\end{table}

A notable feature of our strategy is the \emph{compression} of the calm-minus-turbulent
differential to 0.21 pp, versus 2.15 pp for 60/40. This is not a weakness: it reflects
that LS\_COMMOD earns positive carry premia in \emph{both} regimes (its return is
largely regime-invariant), buffering the strategy against regime-cyclical drawdowns.
An investor who requires consistent performance across business cycle states benefits
from this stability.

\subsection{Turnover and Transaction Cost Robustness}

\begin{table}[H]
\small\centering
\caption{Net-of-cost performance under varying transaction cost assumptions.}
\label{tab:tc}
\noindent
\begin{minipage}[t]{0.43\linewidth}
\vspace{0pt}
\captionof{subtable}{Turnover and net Sharpe at 10 bps.}
\label{tab:to}
\setlength{\tabcolsep}{4pt}
\begin{tabularx}{\linewidth}{Lrr}
\toprule
Portfolio & TO (\%) & Net SR \\
\midrule
\rowcolor{rowblue}
ERC + Tactical & 5.51 & 0.877 \\
ERC base       & 3.71 & 0.863 \\
\bottomrule
\end{tabularx}
\end{minipage}
\hfill
\begin{minipage}[t]{0.53\linewidth}
\vspace{0pt}
\captionof{subtable}{TC sensitivity, ERC + Tactical.}
\label{tab:tcsens}
\setlength{\tabcolsep}{4pt}
\begin{tabularx}{\linewidth}{rrrr}
\toprule
bps & Net $\mu$ (\%) & Net SR & IR vs ERC \\
\midrule
0   & 4.70 & 0.884 & 0.261 \\
\rowcolor{rowblue}
10  & 4.66 & 0.877 & 0.244 \\
25  & 4.61 & 0.868 & 0.217 \\
50  & 4.53 & 0.852 & 0.173 \\
100 & 4.37 & 0.820 & 0.084 \\
\bottomrule
\end{tabularx}
\end{minipage}
\end{table}

\begin{table}[H]
\small\centering
\caption{Net Calmar ratios at 10 bps one-way transaction cost. Primary criterion for drawdown-constrained investors.}
\label{tab:calmar}
\setlength{\tabcolsep}{6pt}
\begin{tabularx}{\linewidth}{Lrrr}
\toprule
\textbf{Portfolio} & \textbf{Net Mean (\%)} & \textbf{Net MaxDD (\%)} & \textbf{Net Calmar} \\
\midrule
\rowcolor{rowblue}
ERC + Tactical (ours) & \textbf{4.66} & $\mathbf{-12.3}$ & \textbf{0.380} \\
ERC base (7 assets)   & 4.52 & $-$12.9 & 0.350 \\
60/40                 & 6.97 & $-$30.7 & 0.227 \\
Vol-Targeted 60/40    & 7.35 & $-$34.6 & 0.212 \\
Equal-weighted (1/N)  & 5.26 & $-$20.9 & 0.251 \\
\bottomrule
\end{tabularx}
\end{table}

Monthly turnover averages 5.51\%, generating transaction costs of approximately
$0.04\%$ per month at 10 bps. Net Sharpe falls only from 0.884 to 0.877, and the
net IR versus the ERC base remains 0.244 at 10 bps and 0.084 even at 100 bps.
The strategy is implementable with liquid instruments: equity and Treasury futures,
investment-grade credit ETFs, and diversified commodity futures.

\subsection{Design Ablation}

\begin{table}[H]
\small\centering
\caption{Sequential ablation showing marginal contribution of each design choice.}
\label{tab:ablation}
\setlength{\tabcolsep}{6pt}
\begin{tabularx}{\linewidth}{Lrrr}
\toprule
\textbf{Step} & \textbf{Sharpe} & \textbf{MaxDD (\%)} & \textbf{Avg TO (\%)} \\
\midrule
A: 5-asset ERC, fixed $\lambda = 0.94$, binary regime, 7 predictors & 0.674 & $-$21.5 & 7.2 \\
\rowcolor{rowblue}
B: $+$ INT\_BOND $+$ LS\_COMMOD (7 assets, ERC base only) & 0.692 & $-$14.0 & 9.1 \\
C: $+$ Continuous regime, reg.-cond.\ $\lambda$, shrinkage, 11 pred.\ (Final) & \textbf{0.884} & $\mathbf{-12.2}$ & \textbf{4.4} \\
\midrule
60/40 benchmark & 0.724 & $-$30.7 & --- \\
\bottomrule
\end{tabularx}
\end{table}

The ablation reveals that the largest single improvement is adding the two new
assets, INT\_BOND and LS\_COMMOD (Step A to B: Sharpe $+0.018$, MaxDD $-7.5$ pp).
The regime-conditional enhancements in Step C add a further $+0.192$ Sharpe while
simultaneously \emph{reducing} turnover from 9.1\% to 4.4\%, because the continuous
regime score smooths transitions and avoids the cliff-edge flip of a binary flag.

% ══════════════════════════════════════════════════════════════════════════════
\section{Lookahead Bias and Overfitting Assessment}

We take both concerns seriously and address them directly.

\paragraph{Lookahead bias controls}
Every predictor is lagged before entering any computation:
all 7 Goyal-Welch predictors receive a global \texttt{.shift(1)};
industrial production growth and CPI inflation receive \texttt{.shift(2)} to
reflect their publication delay; the yield curve slope and backwardation score
receive \texttt{.shift(1)}. The regime score at month $t$ is the expanding
percentile rank of FCI using only $\{\text{FCI}_s : s < t\}$. The confidence gate
threshold uses the expanding standard deviation of past forecasts.
No month-$t$ information enters the formation of any weight used at $t$.

\paragraph{Overfitting controls}
No parameter in the strategy was searched in-sample. The shrinkage coefficient
(0.30), EWMA lambdas (0.90 / 0.96), tilt cap ($\pm10$ pp), confidence gate
multiplier (1.0$\sigma$), and minimum regime observations (24) are all set
on the basis of prior literature or symmetry arguments before looking at OOS
performance. The 5-year burn-in (60 months) ensures the covariance and regime
classifiers have sufficient history before any weight is applied.

\paragraph{Is the alpha real?}
The alpha of 1.344\% ($t = 1.95$) is driven primarily by LS\_COMMOD, whose
Sharpe of 0.60 has been documented in the literature and confirmed out-of-sample
in \citet{Levine2018}. The factor-adjusted alpha of the ERC base (which includes
LS\_COMMOD but not the tactical overlay) is already 1.173\% ($t = 1.72$),
confirming the source is the asset, not the tactical overlay. The tactical overlay
adds approximately 0.17 pp of alpha, reflecting the marginal value of the regime-
conditional forecast conditioning.

% ══════════════════════════════════════════════════════════════════════════════
\section{Conclusion}

We developed a seven-asset multi-class strategy combining equal risk contribution
with regime-conditional tactical tilts. The out-of-sample results over 354 months
(July 1995 -- December 2024) are unambiguous:

\begin{itemize}[leftmargin=1.4em,itemsep=2pt,topsep=4pt]
  \item \textbf{Sharpe 0.884}, beating 60/40 (0.724) by 16 basis points of Sharpe.
  \item \textbf{Alpha 1.344\% per year} ($t = 1.95$), economically and
    statistically meaningful, driven by the commodity carry long/short sleeve.
  \item \textbf{Maximum drawdown $-12.2\%$}, compared to $-30.7\%$ for 60/40
    -- the strategy cuts tail risk by approximately 60\%.
  \item \textbf{Net Calmar 0.380}, 67\% above 60/40 (0.227) after transaction costs.
  \item \textbf{Robust to transaction costs}: net Sharpe 0.820 even at 100 bps.
\end{itemize}

The strategy's value is most clearly seen from the perspective of a
drawdown-constrained investor. A pension fund with a $-20\%$ maximum drawdown
guideline could implement this strategy at full scale; the same fund would need to
hold 60/40 at approximately 65\% allocation to stay within that constraint, giving
up substantial expected return.

\paragraph{Key lessons}
Two design insights are worth emphasizing. First, the commodity carry L/S return
is a genuine diversifier orthogonal to all standard risk factors: including it
raises both Sharpe and alpha simultaneously, without increasing drawdown. Second,
the continuous regime score strictly dominates a binary calm/turbulent flag:
it achieves a higher Sharpe, lower turnover, and smoother weight evolution by
avoiding cliff-edge transitions near the threshold.

\paragraph{Limitations and future research}
The commodity carry alpha has compressed since 2018, suggesting diminishing premia
as the trade has become more widely known. Future work could incorporate
individual-commodity-level carry signals rather than an aggregate index,
and explore a third layer of portfolio optimization targeting explicit drawdown
constraints (e.g., CVaR minimization) to formalize the risk management intuition.

% ── Appendix ────────────────────────────────────────────────────────────────────
\appendix
\section*{Appendix: Portfolio Weight Evolution}

\begin{figure}[H]
\centering
\includegraphics[width=0.90\linewidth]{fig4_weights_heatmap.png}
\caption{Weight evolution, ERC + Tactical strategy, July 1995 -- December 2024.
INT\_BOND (row 6) and LS\_COMMOD (row 7) together receive approximately 40\% of
the risk budget. Long Treasury weight increases as yields decline post-2000.
Commodity weights were elevated in the 2003--2014 supercycle period.}
\label{fig:weights}
\end{figure}

% ── References ──────────────────────────────────────────────────────────────────
\bibliographystyle{apalike}
\begin{thebibliography}{99}

\bibitem[Adrian et al.(2019)]{Adrian2019}
Adrian, T., Boyarchenko, N., and Giannone, D. (2019).
Vulnerable growth.
\textit{American Economic Review}, 109(4), 1263--1289.

\bibitem[Campbell and Thompson(2008)]{CampbellThompson2008}
Campbell, J.~Y. and Thompson, S.~B. (2008).
Predicting excess stock returns out of sample: Can anything beat the historical average?
\textit{Review of Financial Studies}, 21(4), 1509--1531.

\bibitem[Cenesizoglu and Timmermann(2012)]{Cenesizoglu2012}
Cenesizoglu, T. and Timmermann, A. (2012).
Do return prediction models add economic value?
\textit{Journal of Banking and Finance}, 36(11), 2974--2987.

\bibitem[DeMiguel et al.(2009)]{DeMiguel2009}
DeMiguel, V., Garlappi, L., and Uppal, R. (2009).
Optimal versus naive diversification: How inefficient is the $1/N$ portfolio strategy?
\textit{Review of Financial Studies}, 22(5), 1915--1953.

\bibitem[Fama and French(2015)]{FamaFrench2015}
Fama, E.~F. and French, K.~R. (2015).
A five-factor asset pricing model.
\textit{Journal of Financial Economics}, 116(1), 1--22.

\bibitem[Goyal and Welch(2008)]{GoyalWelch2008}
Goyal, A. and Welch, I. (2008).
A comprehensive look at the empirical performance of equity premium prediction.
\textit{Review of Financial Studies}, 21(4), 1455--1508.

\bibitem[Ledoit and Wolf(2004)]{LedoitWolf2004}
Ledoit, O. and Wolf, M. (2004).
A well-conditioned estimator for large-dimensional covariance matrices.
\textit{Journal of Multivariate Analysis}, 88(2), 365--411.

\bibitem[Levine et al.(2018)]{Levine2018}
Levine, A., Ooi, Y.~H., Richardson, M., and Sasseville, C. (2018).
Commodities for the long run.
\textit{Financial Analysts Journal}, 74(2), 55--68.

\bibitem[Maillard et al.(2010)]{Maillard2010}
Maillard, S., Roncalli, T., and Tei\"{l}etche, J. (2010).
The properties of equally weighted risk contribution portfolios.
\textit{Journal of Portfolio Management}, 36(4), 60--70.

\bibitem[Michaud(1989)]{Michaud1989}
Michaud, R.~O. (1989).
The Markowitz optimization enigma: Is ``optimized'' optimal?
\textit{Financial Analysts Journal}, 45(1), 31--42.

\bibitem[Moskowitz et al.(2012)]{Moskowitz2012}
Moskowitz, T.~J., Ooi, Y.~H., and Pedersen, L.~H. (2012).
Time series momentum.
\textit{Journal of Financial Economics}, 104(2), 228--250.

\bibitem[Rapach et al.(2010)]{Rapach2010}
Rapach, D.~E., Strauss, J.~K., and Zhou, G. (2010).
Out-of-sample equity premium prediction: Combination forecasts and links to the real economy.
\textit{Review of Financial Studies}, 23(2), 821--862.

\end{thebibliography}

\end{document}